\subsection{章末確認問題}

\begin{enumerate}
\item 「利用者は注文番号で注文履歴を検索できる」は、どの分類の要求か。
\item 「検索結果は通常時2秒以内に表示される」は、機能要求か非機能要求か。理由も述べよ。
\item 「実装言語はPythonとする」は、技術制約かサービス品質制約か。
\item 要求獲得と要求分析の違いを、自分の言葉で説明せよ。
\item 要求仕様化において、自然言語、受け入れ基準、モデルの違いを説明せよ。
\item 要求妥当性確認で、レビュー、シミュレーション、プロトタイプはどのように使い分けるか。
\item 要求変更を受け入れるとき、なぜ費用・納期・スコープへの影響を確認する必要があるか。
\item 追跡可能性が、変更時の影響分析に役立つ理由を説明せよ。
\end{enumerate}

\MiniBox{解答の方向性}{1は機能要求。2はサービス品質制約であり非機能要求。3は技術制約。4以降は、本文の「合意を作る活動」「候補要求を表面化させる」「意味・矛盾・実現可能性を調べる」「合意を維持する」といった語を使って説明できる。}
